\section{Experiments}
\label{sec:experiments}

\subsection{Setup}
\label{sec:setup}

A model cannot be separated from the framework that runs it, so what is under test here is not a
model but the whole combination of model, harness and reasoning effort, which we call a
\emph{system}. We evaluate six systems: three GPT-5.6 variants --- Sol, Terra and Luna --- under
Codex at all six effort levels; two Claude~5 variants, Opus~5 and Sonnet~5, under Claude Code at the
five levels that harness exposes; and Kimi~K3 under Claude Code at its highest. That is 29
configurations, each attempting all ten tasks, for 290 cells.

\subsection{Results}
\label{sec:results}

% All figures in this subsection are read off Tables~\ref{tab:grid} and
% \ref{tab:mapped} and computed in build/lark/mapped_table.csv.
\textbf{The whole study sits in the lowest fifth of the scale.} The mean score over the 290 cells is
$0.166$, the strongest system averages $0.250$ (Figure~\ref{fig:bysystem}), and the single best
configuration in the study, Claude Opus~5 at medium effort, averages $0.288$. Against a scale on
which $0.1$ is the algorithm each repository already ships and $1.0$ is the task optimum, that is a
fifth of the distance at the very top and well under a tenth on average. In the other direction, 124
of the 290 cells fall below $0.1$: more than two fifths of the attempts leave the repository with
something worse than what it had. None of this is legible in the raw units of
Table~\ref{tab:grid}, where lifting weight averaging by $0.015$ of top-1 accuracy and taking a
perplexity from $53.4$ to $13.0$ are both simply cells that beat a baseline.

\textbf{Systems are ordered, and the ordering is compressed.} Figure~\ref{fig:bysystem} separates the
six cleanly --- Claude Opus~5 at $0.250$, then GPT-5.6 Sol at $0.191$, Kimi~K3 at $0.174$, Claude
Sonnet~5 at $0.145$, GPT-5.6 Terra at $0.135$ and GPT-5.6 Luna at $0.117$. But the entire range lies inside the bottom quarter of the scale, so the choice of
system moves the number without moving the regime: the best system's average is closer to the
weakest system's than it is to the optimum it was asked to approach.

\textbf{Spend does not explain the result.} The cost column of Table~\ref{tab:grid} ranges about
ninefold across systems on the same harness --- a median configuration costs \$434 for Sol and \$48
for Luna --- and the ordering it induces is not the ordering of the scores. Opus~5 leads the study at
a median of \$181, under half of what the second-placed system spent, and Sonnet~5 spends about twice
Luna's budget for a $0.028$ difference. Whatever is separating these systems, it is not how much
exploration they bought.

% Generated by scripts/make_grid_table.py -- do not edit by hand.
\begin{table}[t]
\centering
\caption{\textbf{Every configuration on every task, in the metric's own units.} Rows are the 29 (model, harness, effort) configurations; \textbf{\$} is what the four hours of exploration cost. \emph{Baseline} is the score of the repository's own code under the identical procedure (\S\ref{sec:baselines}). Backgrounds report the mapped score of \S\ref{sec:scoring}: \colorbox{bandlow}{below $0.1$} is worse than that baseline, \colorbox{bandmid}{$0.1$ to $0.4$} beats it, \colorbox{bandhigh}{above $0.4$} closes more than a third of the distance to the optimum. \textbf{Bold} marks the best cell in a column, and a dash a configuration that returned nothing trainable. The last three rows give the ladder each column is scored on.}
\label{tab:grid}
\scriptsize
\begin{adjustbox}{max width=\linewidth}
\begin{tabular}{llllrrrrrrrrrrr}
\toprule
& \textbf{System} & \textbf{Harness} & \textbf{Effort} & \textbf{OpenR1}\,$\uparrow$ & \textbf{RAGEN}\,$\uparrow$ & \textbf{OPD}\,$\uparrow$ & \textbf{BTRM}\,$\uparrow$ & \textbf{DPO}\,$\uparrow$ & \textbf{DDPO}\,$\uparrow$ & \textbf{NPO}\,$\uparrow$ & \textbf{DiGress}\,$\downarrow$ & \textbf{Soup}\,$\uparrow$ & \textbf{OWL}\,$\downarrow$ & \textbf{Cost (USD)}\\
\midrule
& \emph{Baseline} & \na{} & \na{} & 0.127 & 0.170 & 0.436 & 74.9 & 0.424 & 5.84 & 0.887 & 65.8 & 0.686 & 53.4 & \na{} \\
\midrule
\raisebox{-0.15ex}{
\includegraphics[height=1.5ex]{logos/claude}} & Claude Opus 5 & Claude Code & \texttt{low} & \cellcolor{bandmid}\textbf{0.138} & \cellcolor{bandhigh}1.00 & \cellcolor{bandlow}0.421 & \cellcolor{bandlow}71.6 & \cellcolor{bandmid}0.467 & \cellcolor{bandhigh}12.1 & \cellcolor{bandmid}0.997 & \cellcolor{bandmid}65.4 & \cellcolor{bandmid}0.696 & \cellcolor{bandhigh}\textbf{13.0} & 181 \\
\raisebox{-0.15ex}{
\includegraphics[height=1.5ex]{logos/claude}} & Claude Opus 5 & Claude Code & \texttt{medium} & \cellcolor{bandlow}0.121 & \cellcolor{bandhigh}1.00 & \cellcolor{bandlow}0.432 & \cellcolor{bandmid}\textbf{77.1} & \cellcolor{bandhigh}\textbf{0.622} & \cellcolor{bandmid}8.98 & \cellcolor{bandmid}1.01 & \cellcolor{bandmid}65.8 & \cellcolor{bandmid}0.696 & \cellcolor{bandhigh}13.4 & 166 \\
\raisebox{-0.15ex}{
\includegraphics[height=1.5ex]{logos/claude}} & Claude Opus 5 & Claude Code & \texttt{high} & \cellcolor{bandmid}0.127 & \cellcolor{bandhigh}1.00 & \cellcolor{bandmid}\textbf{0.449} & \cellcolor{bandlow}71.9 & \cellcolor{bandmid}0.615 & \cellcolor{bandmid}8.42 & \cellcolor{bandmid}1.01 & \cellcolor{bandmid}65.7 & \cellcolor{bandmid}0.694 & \cellcolor{bandhigh}13.2 & 181 \\
\raisebox{-0.15ex}{
\includegraphics[height=1.5ex]{logos/claude}} & Claude Opus 5 & Claude Code & \texttt{xhigh} & \cellcolor{bandmid}0.128 & \cellcolor{bandmid}0.211 & \cellcolor{bandmid}0.440 & \cellcolor{bandlow}64.4 & \cellcolor{bandlow}0.421 & \cellcolor{bandhigh}14.4 & \cellcolor{bandmid}1.02 & \cellcolor{bandmid}65.8 & \cellcolor{bandmid}0.696 & \cellcolor{bandhigh}13.3 & 185 \\
\raisebox{-0.15ex}{
\includegraphics[height=1.5ex]{logos/claude}} & Claude Opus 5 & Claude Code & \texttt{max} & \cellcolor{bandlow}0.125 & \cellcolor{bandmid}0.234 & \cellcolor{bandlow}0.392 & \cellcolor{bandmid}75.7 & \cellcolor{bandmid}0.477 & \cellcolor{bandhigh}\textbf{17.7} & \cellcolor{bandmid}\textbf{1.03} & \cellcolor{bandmid}64.8 & \cellcolor{bandmid}0.694 & \cellcolor{bandhigh}13.0 & 195 \\
\addlinespace[2pt]
\raisebox{-0.15ex}{
\includegraphics[height=1.5ex]{logos/claude}} & Claude Sonnet 5 & Claude Code & \texttt{low} & \cellcolor{bandlow}0.108 & \cellcolor{bandlow}0.043 & \cellcolor{bandlow}0.432 & \cellcolor{bandlow}69.3 & \cellcolor{bandlow}0.404 & \cellcolor{bandmid}6.13 & \cellcolor{bandmid}0.969 & \cellcolor{bandlow}66.3 & \cellcolor{bandmid}0.698 & \cellcolor{bandlow}54.0 & 93 \\
\raisebox{-0.15ex}{
\includegraphics[height=1.5ex]{logos/claude}} & Claude Sonnet 5 & Claude Code & \texttt{medium} & \cellcolor{bandlow}-- & \cellcolor{bandlow}0.000 & \cellcolor{bandlow}0.427 & \cellcolor{bandlow}71.5 & \cellcolor{bandlow}0.400 & \cellcolor{bandmid}5.86 & \cellcolor{bandmid}0.969 & \cellcolor{bandmid}65.4 & \cellcolor{bandmid}0.693 & \cellcolor{bandmid}19.7 & 98 \\
\raisebox{-0.15ex}{
\includegraphics[height=1.5ex]{logos/claude}} & Claude Sonnet 5 & Claude Code & \texttt{high} & \cellcolor{bandlow}0.113 & \cellcolor{bandhigh}0.895 & \cellcolor{bandlow}0.427 & \cellcolor{bandlow}74.5 & \cellcolor{bandmid}0.540 & \cellcolor{bandlow}5.82 & \cellcolor{bandlow}0.733 & \cellcolor{bandlow}66.5 & \cellcolor{bandmid}0.694 & \cellcolor{bandmid}21.7 & 96 \\
\raisebox{-0.15ex}{
\includegraphics[height=1.5ex]{logos/claude}} & Claude Sonnet 5 & Claude Code & \texttt{xhigh} & \cellcolor{bandlow}0.114 & \cellcolor{bandmid}0.240 & \cellcolor{bandlow}0.421 & \cellcolor{bandlow}38.3 & \cellcolor{bandmid}0.506 & \cellcolor{bandlow}4.49 & \cellcolor{bandmid}0.963 & \cellcolor{bandlow}66.6 & \cellcolor{bandmid}0.694 & \cellcolor{bandmid}21.2 & 101 \\
\raisebox{-0.15ex}{
\includegraphics[height=1.5ex]{logos/claude}} & Claude Sonnet 5 & Claude Code & \texttt{max} & \cellcolor{bandlow}0.113 & \cellcolor{bandmid}0.232 & \cellcolor{bandlow}0.428 & \cellcolor{bandmid}74.9 & \cellcolor{bandmid}0.542 & \cellcolor{bandmid}6.15 & \cellcolor{bandmid}0.948 & \cellcolor{bandlow}66.2 & \cellcolor{bandmid}0.699 & \cellcolor{bandmid}14.3 & 109 \\
\addlinespace[2pt]
\raisebox{-0.15ex}{
\includegraphics[height=1.5ex]{logos/kimi}} & Kimi K3 & Claude Code & \texttt{max} & \cellcolor{bandlow}0.099 & \cellcolor{bandmid}0.186 & \cellcolor{bandlow}0.405 & \cellcolor{bandlow}74.1 & \cellcolor{bandmid}0.542 & \cellcolor{bandmid}9.04 & \cellcolor{bandmid}0.986 & \cellcolor{bandmid}65.2 & \cellcolor{bandmid}0.694 & \cellcolor{bandmid}14.3 & 30 \\
\addlinespace[2pt]
\raisebox{-0.15ex}{
\includegraphics[height=1.5ex]{logos/openai}} & GPT-5.6 Sol & Codex & \texttt{none} & \cellcolor{bandlow}0.115 & \cellcolor{bandlow}0.092 & \cellcolor{bandlow}0.426 & \cellcolor{bandmid}75.2 & \cellcolor{bandlow}0.400 & \cellcolor{bandlow}5.82 & \cellcolor{bandmid}0.994 & \cellcolor{bandlow}67.0 & \cellcolor{bandmid}0.694 & \cellcolor{bandmid}20.9 & 339 \\
\raisebox{-0.15ex}{
\includegraphics[height=1.5ex]{logos/openai}} & GPT-5.6 Sol & Codex & \texttt{low} & \cellcolor{bandlow}0.127 & \cellcolor{bandhigh}0.883 & \cellcolor{bandmid}0.437 & \cellcolor{bandlow}74.0 & \cellcolor{bandmid}0.431 & \cellcolor{bandlow}-- & \cellcolor{bandlow}0.746 & \cellcolor{bandlow}66.9 & \cellcolor{bandmid}0.694 & \cellcolor{bandmid}14.9 & 337 \\
\raisebox{-0.15ex}{
\includegraphics[height=1.5ex]{logos/openai}} & GPT-5.6 Sol & Codex & \texttt{medium} & \cellcolor{bandlow}0.124 & \cellcolor{bandlow}-- & \cellcolor{bandlow}0.430 & \cellcolor{bandlow}73.7 & \cellcolor{bandmid}0.494 & \cellcolor{bandmid}6.83 & \cellcolor{bandmid}0.998 & \cellcolor{bandmid}65.5 & \cellcolor{bandmid}0.695 & \cellcolor{bandmid}15.8 & 449 \\
\raisebox{-0.15ex}{
\includegraphics[height=1.5ex]{logos/openai}} & GPT-5.6 Sol & Codex & \texttt{high} & \cellcolor{bandlow}0.126 & \cellcolor{bandhigh}1.00 & \cellcolor{bandlow}0.424 & \cellcolor{bandlow}74.0 & \cellcolor{bandmid}0.458 & \cellcolor{bandlow}5.77 & \cellcolor{bandmid}0.966 & \cellcolor{bandlow}66.0 & \cellcolor{bandmid}0.697 & \cellcolor{bandhigh}13.3 & 420 \\
\raisebox{-0.15ex}{
\includegraphics[height=1.5ex]{logos/openai}} & GPT-5.6 Sol & Codex & \texttt{xhigh} & \cellcolor{bandlow}0.125 & \cellcolor{bandhigh}0.549 & \cellcolor{bandlow}0.418 & \cellcolor{bandlow}73.8 & \cellcolor{bandmid}0.482 & \cellcolor{bandmid}5.92 & \cellcolor{bandmid}1.03 & \cellcolor{bandmid}\textbf{63.7} & \cellcolor{bandmid}0.692 & \cellcolor{bandhigh}13.9 & 521 \\
\raisebox{-0.15ex}{
\includegraphics[height=1.5ex]{logos/openai}} & GPT-5.6 Sol & Codex & \texttt{max} & \cellcolor{bandlow}0.126 & \cellcolor{bandhigh}\textbf{1.00} & \cellcolor{bandlow}0.429 & \cellcolor{bandlow}73.9 & \cellcolor{bandmid}0.436 & \cellcolor{bandmid}8.98 & \cellcolor{bandmid}0.998 & \cellcolor{bandlow}65.9 & \cellcolor{bandmid}\textbf{0.701} & \cellcolor{bandmid}17.4 & 626 \\
\addlinespace[2pt]
\raisebox{-0.15ex}{
\includegraphics[height=1.5ex]{logos/openai}} & GPT-5.6 Terra & Codex & \texttt{none} & \cellcolor{bandlow}0.096 & \cellcolor{bandmid}0.221 & \cellcolor{bandlow}-- & \cellcolor{bandmid}75.2 & \cellcolor{bandmid}0.443 & \cellcolor{bandlow}5.54 & \cellcolor{bandlow}0.713 & \cellcolor{bandlow}-- & \cellcolor{bandlow}-- & \cellcolor{bandlow}-- & 4 \\
\raisebox{-0.15ex}{
\includegraphics[height=1.5ex]{logos/openai}} & GPT-5.6 Terra & Codex & \texttt{low} & \cellcolor{bandlow}0.127 & \cellcolor{bandlow}-- & \cellcolor{bandlow}-- & \cellcolor{bandlow}72.9 & \cellcolor{bandlow}0.421 & \cellcolor{bandlow}5.54 & \cellcolor{bandmid}0.952 & \cellcolor{bandmid}65.3 & \cellcolor{bandmid}0.694 & \cellcolor{bandmid}20.4 & 35 \\
\raisebox{-0.15ex}{
\includegraphics[height=1.5ex]{logos/openai}} & GPT-5.6 Terra & Codex & \texttt{medium} & \cellcolor{bandlow}0.126 & \cellcolor{bandmid}0.184 & \cellcolor{bandmid}0.443 & \cellcolor{bandlow}74.1 & \cellcolor{bandmid}0.448 & \cellcolor{bandmid}5.85 & \cellcolor{bandmid}0.921 & \cellcolor{bandmid}64.3 & \cellcolor{bandmid}0.694 & \cellcolor{bandmid}20.2 & 43 \\
\raisebox{-0.15ex}{
\includegraphics[height=1.5ex]{logos/openai}} & GPT-5.6 Terra & Codex & \texttt{high} & \cellcolor{bandlow}0.108 & \cellcolor{bandlow}0.086 & \cellcolor{bandlow}-- & \cellcolor{bandmid}75.4 & \cellcolor{bandmid}0.494 & \cellcolor{bandmid}6.62 & \cellcolor{bandmid}0.997 & \cellcolor{bandlow}-- & \cellcolor{bandmid}0.694 & \cellcolor{bandmid}15.2 & 214 \\
\raisebox{-0.15ex}{
\includegraphics[height=1.5ex]{logos/openai}} & GPT-5.6 Terra & Codex & \texttt{xhigh} & \cellcolor{bandmid}0.127 & \cellcolor{bandmid}0.266 & \cellcolor{bandlow}0.433 & \cellcolor{bandlow}70.9 & \cellcolor{bandmid}0.438 & \cellcolor{bandmid}5.96 & \cellcolor{bandmid}0.923 & \cellcolor{bandmid}65.6 & \cellcolor{bandmid}0.694 & \cellcolor{bandmid}16.2 & 229 \\
\raisebox{-0.15ex}{
\includegraphics[height=1.5ex]{logos/openai}} & GPT-5.6 Terra & Codex & \texttt{max} & \cellcolor{bandlow}0.127 & \cellcolor{bandhigh}0.998 & \cellcolor{bandlow}0.435 & \cellcolor{bandlow}66.9 & \cellcolor{bandmid}0.460 & \cellcolor{bandmid}6.01 & \cellcolor{bandmid}0.954 & \cellcolor{bandlow}69.9 & \cellcolor{bandmid}0.694 & \cellcolor{bandmid}15.2 & 346 \\
\addlinespace[2pt]
\raisebox{-0.15ex}{
\includegraphics[height=1.5ex]{logos/openai}} & GPT-5.6 Luna & Codex & \texttt{none} & \cellcolor{bandlow}0.126 & \cellcolor{bandlow}-- & \cellcolor{bandlow}0.432 & \cellcolor{bandmid}75.3 & \cellcolor{bandlow}0.407 & \cellcolor{bandlow}5.11 & \cellcolor{bandmid}0.936 & \cellcolor{bandlow}-- & \cellcolor{bandmid}0.694 & \cellcolor{bandmid}42.9 & 17 \\
\raisebox{-0.15ex}{
\includegraphics[height=1.5ex]{logos/openai}} & GPT-5.6 Luna & Codex & \texttt{low} & \cellcolor{bandlow}0.115 & \cellcolor{bandlow}0.154 & \cellcolor{bandlow}-- & \cellcolor{bandlow}74.5 & \cellcolor{bandmid}0.472 & \cellcolor{bandlow}5.54 & \cellcolor{bandlow}-- & \cellcolor{bandlow}-- & \cellcolor{bandmid}0.694 & \cellcolor{bandmid}51.0 & 7 \\
\raisebox{-0.15ex}{
\includegraphics[height=1.5ex]{logos/openai}} & GPT-5.6 Luna & Codex & \texttt{medium} & \cellcolor{bandlow}0.125 & \cellcolor{bandlow}-- & \cellcolor{bandmid}0.438 & \cellcolor{bandmid}75.8 & \cellcolor{bandmid}0.492 & \cellcolor{bandlow}5.56 & \cellcolor{bandlow}0.748 & \cellcolor{bandlow}66.2 & \cellcolor{bandmid}0.694 & \cellcolor{bandmid}22.4 & 30 \\
\raisebox{-0.15ex}{
\includegraphics[height=1.5ex]{logos/openai}} & GPT-5.6 Luna & Codex & \texttt{high} & \cellcolor{bandlow}0.122 & \cellcolor{bandlow}0.125 & \cellcolor{bandlow}0.424 & \cellcolor{bandlow}74.0 & \cellcolor{bandlow}0.402 & \cellcolor{bandlow}5.62 & \cellcolor{bandmid}0.990 & \cellcolor{bandlow}67.4 & \cellcolor{bandmid}0.694 & \cellcolor{bandmid}29.2 & 66 \\
\raisebox{-0.15ex}{
\includegraphics[height=1.5ex]{logos/openai}} & GPT-5.6 Luna & Codex & \texttt{xhigh} & \cellcolor{bandlow}0.126 & \cellcolor{bandlow}-- & \cellcolor{bandlow}0.427 & \cellcolor{bandmid}75.4 & \cellcolor{bandmid}0.564 & \cellcolor{bandlow}5.74 & \cellcolor{bandmid}0.957 & \cellcolor{bandlow}68.3 & \cellcolor{bandmid}0.692 & \cellcolor{bandmid}16.2 & 110 \\
\raisebox{-0.15ex}{
\includegraphics[height=1.5ex]{logos/openai}} & GPT-5.6 Luna & Codex & \texttt{max} & \cellcolor{bandlow}0.116 & \cellcolor{bandlow}-- & \cellcolor{bandlow}0.433 & \cellcolor{bandlow}74.4 & \cellcolor{bandmid}0.533 & \cellcolor{bandmid}6.53 & \cellcolor{bandmid}0.945 & \cellcolor{bandmid}65.3 & \cellcolor{bandmid}0.694 & \cellcolor{bandmid}19.7 & 108 \\
\midrule
\multicolumn{4}{l}{\emph{Optimum} $x^{\ast}$} & 1 & 1 & 1 & 100 & 1 & 23.23 & 2.08 & 0 & 1 & 1 & \na{} \\
\multicolumn{4}{l}{\emph{Uninformative} $x_{\perp}$} & 0 & 0 & 0 & 50 & 0 & $-12.76$ & 0 & $\infty$ & 0 & $\infty$ & \na{} \\
\multicolumn{4}{l}{\emph{Coordinate} $\varphi$} & $x$ & $x$ & $x$ & $x$ & $x$ & $x$ & $x$ & $-x$ & $x$ & $-\log x$ & \na{} \\
\bottomrule
\end{tabular}
\end{adjustbox}
\end{table}


% Generated by scripts/make_grid_table.py -- do not edit by hand.
\begin{table}[t]
\centering
\caption{\textbf{The same cells after the ladder of \S\ref{sec:scoring}.} $0.1$ is the repository's own recipe and $1.0$ the task optimum, so a score states how much of the remaining distance a submission closed; a configuration that returned nothing scores $0$. Backgrounds and bold follow Table~\ref{tab:grid}.}
\label{tab:mapped}
\scriptsize
\begin{adjustbox}{max width=\linewidth}
\begin{tabular}{llllrrrrrrrrrrr}
\toprule
& \textbf{System} & \textbf{Harness} & \textbf{Effort} & \textbf{OpenR1} & \textbf{RAGEN} & \textbf{OPD} & \textbf{BTRM} & \textbf{DPO} & \textbf{DDPO} & \textbf{NPO} & \textbf{DiGress} & \textbf{Soup} & \textbf{OWL} & \textbf{avg}\\
\midrule
\raisebox{-0.15ex}{
\includegraphics[height=1.5ex]{logos/claude}} & Claude Opus 5 & Claude Code & \texttt{low} & \cellcolor{bandmid}\textbf{0.111} & \cellcolor{bandhigh}1.000 & \cellcolor{bandlow}0.097 & \cellcolor{bandlow}0.087 & \cellcolor{bandmid}0.167 & \cellcolor{bandhigh}0.424 & \cellcolor{bandmid}0.182 & \cellcolor{bandmid}0.105 & \cellcolor{bandmid}0.129 & \cellcolor{bandhigh}\textbf{0.420} & 0.272 \\
\raisebox{-0.15ex}{
\includegraphics[height=1.5ex]{logos/claude}} & Claude Opus 5 & Claude Code & \texttt{medium} & \cellcolor{bandlow}0.095 & \cellcolor{bandhigh}1.000 & \cellcolor{bandlow}0.099 & \cellcolor{bandmid}\textbf{0.179} & \cellcolor{bandhigh}\textbf{0.409} & \cellcolor{bandmid}0.263 & \cellcolor{bandmid}0.195 & \cellcolor{bandmid}0.100 & \cellcolor{bandmid}0.129 & \cellcolor{bandhigh}0.413 & \textbf{0.288} \\
\raisebox{-0.15ex}{
\includegraphics[height=1.5ex]{logos/claude}} & Claude Opus 5 & Claude Code & \texttt{high} & \cellcolor{bandmid}0.100 & \cellcolor{bandhigh}1.000 & \cellcolor{bandmid}\textbf{0.121} & \cellcolor{bandlow}0.088 & \cellcolor{bandmid}0.398 & \cellcolor{bandmid}0.234 & \cellcolor{bandmid}0.195 & \cellcolor{bandmid}0.101 & \cellcolor{bandmid}0.123 & \cellcolor{bandhigh}0.416 & 0.278 \\
\raisebox{-0.15ex}{
\includegraphics[height=1.5ex]{logos/claude}} & Claude Opus 5 & Claude Code & \texttt{xhigh} & \cellcolor{bandmid}0.101 & \cellcolor{bandmid}0.144 & \cellcolor{bandmid}0.106 & \cellcolor{bandlow}0.058 & \cellcolor{bandlow}0.099 & \cellcolor{bandhigh}0.543 & \cellcolor{bandmid}0.203 & \cellcolor{bandmid}0.100 & \cellcolor{bandmid}0.129 & \cellcolor{bandhigh}0.414 & 0.190 \\
\raisebox{-0.15ex}{
\includegraphics[height=1.5ex]{logos/claude}} & Claude Opus 5 & Claude Code & \texttt{max} & \cellcolor{bandlow}0.098 & \cellcolor{bandmid}0.169 & \cellcolor{bandlow}0.090 & \cellcolor{bandmid}0.129 & \cellcolor{bandmid}0.183 & \cellcolor{bandhigh}\textbf{0.714} & \cellcolor{bandmid}\textbf{0.207} & \cellcolor{bandmid}0.114 & \cellcolor{bandmid}0.123 & \cellcolor{bandhigh}0.420 & 0.225 \\
\addlinespace[2pt]
\raisebox{-0.15ex}{
\includegraphics[height=1.5ex]{logos/claude}} & Claude Sonnet 5 & Claude Code & \texttt{low} & \cellcolor{bandlow}0.085 & \cellcolor{bandlow}0.025 & \cellcolor{bandlow}0.099 & \cellcolor{bandlow}0.077 & \cellcolor{bandlow}0.095 & \cellcolor{bandmid}0.115 & \cellcolor{bandmid}0.161 & \cellcolor{bandlow}0.099 & \cellcolor{bandmid}0.134 & \cellcolor{bandlow}0.099 & 0.099 \\
\raisebox{-0.15ex}{
\includegraphics[height=1.5ex]{logos/claude}} & Claude Sonnet 5 & Claude Code & \texttt{medium} & \cellcolor{bandlow}0.000 & \cellcolor{bandlow}0.000 & \cellcolor{bandlow}0.098 & \cellcolor{bandlow}0.086 & \cellcolor{bandlow}0.094 & \cellcolor{bandmid}0.101 & \cellcolor{bandmid}0.162 & \cellcolor{bandmid}0.105 & \cellcolor{bandmid}0.120 & \cellcolor{bandmid}0.326 & 0.109 \\
\raisebox{-0.15ex}{
\includegraphics[height=1.5ex]{logos/claude}} & Claude Sonnet 5 & Claude Code & \texttt{high} & \cellcolor{bandlow}0.088 & \cellcolor{bandhigh}0.886 & \cellcolor{bandlow}0.098 & \cellcolor{bandlow}0.098 & \cellcolor{bandmid}0.281 & \cellcolor{bandlow}0.100 & \cellcolor{bandlow}0.083 & \cellcolor{bandlow}0.099 & \cellcolor{bandmid}0.123 & \cellcolor{bandmid}0.304 & 0.216 \\
\raisebox{-0.15ex}{
\includegraphics[height=1.5ex]{logos/claude}} & Claude Sonnet 5 & Claude Code & \texttt{xhigh} & \cellcolor{bandlow}0.089 & \cellcolor{bandmid}0.176 & \cellcolor{bandlow}0.097 & \cellcolor{bandlow}0.000 & \cellcolor{bandmid}0.228 & \cellcolor{bandlow}0.093 & \cellcolor{bandmid}0.157 & \cellcolor{bandlow}0.099 & \cellcolor{bandmid}0.123 & \cellcolor{bandmid}0.309 & 0.137 \\
\raisebox{-0.15ex}{
\includegraphics[height=1.5ex]{logos/claude}} & Claude Sonnet 5 & Claude Code & \texttt{max} & \cellcolor{bandlow}0.088 & \cellcolor{bandmid}0.167 & \cellcolor{bandlow}0.098 & \cellcolor{bandmid}0.100 & \cellcolor{bandmid}0.284 & \cellcolor{bandmid}0.116 & \cellcolor{bandmid}0.146 & \cellcolor{bandlow}0.099 & \cellcolor{bandmid}0.137 & \cellcolor{bandmid}0.398 & 0.163 \\
\addlinespace[2pt]
\raisebox{-0.15ex}{
\includegraphics[height=1.5ex]{logos/kimi}} & Kimi K3 & Claude Code & \texttt{max} & \cellcolor{bandlow}0.078 & \cellcolor{bandmid}0.117 & \cellcolor{bandlow}0.093 & \cellcolor{bandlow}0.097 & \cellcolor{bandmid}0.284 & \cellcolor{bandmid}0.266 & \cellcolor{bandmid}0.175 & \cellcolor{bandmid}0.108 & \cellcolor{bandmid}0.123 & \cellcolor{bandmid}0.398 & 0.174 \\
\addlinespace[2pt]
\raisebox{-0.15ex}{
\includegraphics[height=1.5ex]{logos/openai}} & GPT-5.6 Sol & Codex & \texttt{none} & \cellcolor{bandlow}0.091 & \cellcolor{bandlow}0.054 & \cellcolor{bandlow}0.098 & \cellcolor{bandmid}0.111 & \cellcolor{bandlow}0.094 & \cellcolor{bandlow}0.100 & \cellcolor{bandmid}0.180 & \cellcolor{bandlow}0.098 & \cellcolor{bandmid}0.123 & \cellcolor{bandmid}0.312 & 0.126 \\
\raisebox{-0.15ex}{
\includegraphics[height=1.5ex]{logos/openai}} & GPT-5.6 Sol & Codex & \texttt{low} & \cellcolor{bandlow}0.100 & \cellcolor{bandhigh}0.873 & \cellcolor{bandmid}0.102 & \cellcolor{bandlow}0.096 & \cellcolor{bandmid}0.111 & \cellcolor{bandlow}0.000 & \cellcolor{bandlow}0.084 & \cellcolor{bandlow}0.098 & \cellcolor{bandmid}0.123 & \cellcolor{bandmid}0.389 & 0.198 \\
\raisebox{-0.15ex}{
\includegraphics[height=1.5ex]{logos/openai}} & GPT-5.6 Sol & Codex & \texttt{medium} & \cellcolor{bandlow}0.097 & \cellcolor{bandlow}0.000 & \cellcolor{bandlow}0.099 & \cellcolor{bandlow}0.095 & \cellcolor{bandmid}0.209 & \cellcolor{bandmid}0.151 & \cellcolor{bandmid}0.183 & \cellcolor{bandmid}0.104 & \cellcolor{bandmid}0.126 & \cellcolor{bandmid}0.376 & 0.144 \\
\raisebox{-0.15ex}{
\includegraphics[height=1.5ex]{logos/openai}} & GPT-5.6 Sol & Codex & \texttt{high} & \cellcolor{bandlow}0.099 & \cellcolor{bandhigh}1.000 & \cellcolor{bandlow}0.097 & \cellcolor{bandlow}0.096 & \cellcolor{bandmid}0.153 & \cellcolor{bandlow}0.100 & \cellcolor{bandmid}0.159 & \cellcolor{bandlow}0.100 & \cellcolor{bandmid}0.132 & \cellcolor{bandhigh}0.414 & 0.235 \\
\raisebox{-0.15ex}{
\includegraphics[height=1.5ex]{logos/openai}} & GPT-5.6 Sol & Codex & \texttt{xhigh} & \cellcolor{bandlow}0.098 & \cellcolor{bandhigh}0.511 & \cellcolor{bandlow}0.096 & \cellcolor{bandlow}0.096 & \cellcolor{bandmid}0.191 & \cellcolor{bandmid}0.104 & \cellcolor{bandmid}0.207 & \cellcolor{bandmid}\textbf{0.129} & \cellcolor{bandmid}0.117 & \cellcolor{bandhigh}0.405 & 0.195 \\
\raisebox{-0.15ex}{
\includegraphics[height=1.5ex]{logos/openai}} & GPT-5.6 Sol & Codex & \texttt{max} & \cellcolor{bandlow}0.099 & \cellcolor{bandhigh}\textbf{1.000} & \cellcolor{bandlow}0.098 & \cellcolor{bandlow}0.096 & \cellcolor{bandmid}0.119 & \cellcolor{bandmid}0.263 & \cellcolor{bandmid}0.183 & \cellcolor{bandlow}0.100 & \cellcolor{bandmid}\textbf{0.143} & \cellcolor{bandmid}0.354 & 0.245 \\
\addlinespace[2pt]
\raisebox{-0.15ex}{
\includegraphics[height=1.5ex]{logos/openai}} & GPT-5.6 Terra & Codex & \texttt{none} & \cellcolor{bandlow}0.075 & \cellcolor{bandmid}0.155 & \cellcolor{bandlow}0.000 & \cellcolor{bandmid}0.111 & \cellcolor{bandmid}0.130 & \cellcolor{bandlow}0.098 & \cellcolor{bandlow}0.080 & \cellcolor{bandlow}0.000 & \cellcolor{bandlow}0.000 & \cellcolor{bandlow}0.000 & 0.065 \\
\raisebox{-0.15ex}{
\includegraphics[height=1.5ex]{logos/openai}} & GPT-5.6 Terra & Codex & \texttt{low} & \cellcolor{bandlow}0.100 & \cellcolor{bandlow}0.000 & \cellcolor{bandlow}0.000 & \cellcolor{bandlow}0.092 & \cellcolor{bandlow}0.099 & \cellcolor{bandlow}0.098 & \cellcolor{bandmid}0.149 & \cellcolor{bandmid}0.107 & \cellcolor{bandmid}0.123 & \cellcolor{bandmid}0.318 & 0.109 \\
\raisebox{-0.15ex}{
\includegraphics[height=1.5ex]{logos/openai}} & GPT-5.6 Terra & Codex & \texttt{medium} & \cellcolor{bandlow}0.099 & \cellcolor{bandmid}0.115 & \cellcolor{bandmid}0.111 & \cellcolor{bandlow}0.097 & \cellcolor{bandmid}0.138 & \cellcolor{bandmid}0.101 & \cellcolor{bandmid}0.125 & \cellcolor{bandmid}0.120 & \cellcolor{bandmid}0.123 & \cellcolor{bandmid}0.320 & 0.135 \\
\raisebox{-0.15ex}{
\includegraphics[height=1.5ex]{logos/openai}} & GPT-5.6 Terra & Codex & \texttt{high} & \cellcolor{bandlow}0.085 & \cellcolor{bandlow}0.051 & \cellcolor{bandlow}0.000 & \cellcolor{bandmid}0.118 & \cellcolor{bandmid}0.209 & \cellcolor{bandmid}0.140 & \cellcolor{bandmid}0.183 & \cellcolor{bandlow}0.000 & \cellcolor{bandmid}0.123 & \cellcolor{bandmid}0.384 & 0.129 \\
\raisebox{-0.15ex}{
\includegraphics[height=1.5ex]{logos/openai}} & GPT-5.6 Terra & Codex & \texttt{xhigh} & \cellcolor{bandmid}0.100 & \cellcolor{bandmid}0.204 & \cellcolor{bandlow}0.099 & \cellcolor{bandlow}0.084 & \cellcolor{bandmid}0.122 & \cellcolor{bandmid}0.106 & \cellcolor{bandmid}0.127 & \cellcolor{bandmid}0.103 & \cellcolor{bandmid}0.123 & \cellcolor{bandmid}0.370 & 0.144 \\
\raisebox{-0.15ex}{
\includegraphics[height=1.5ex]{logos/openai}} & GPT-5.6 Terra & Codex & \texttt{max} & \cellcolor{bandlow}0.100 & \cellcolor{bandhigh}0.998 & \cellcolor{bandlow}0.100 & \cellcolor{bandlow}0.068 & \cellcolor{bandmid}0.156 & \cellcolor{bandmid}0.109 & \cellcolor{bandmid}0.150 & \cellcolor{bandlow}0.094 & \cellcolor{bandmid}0.123 & \cellcolor{bandmid}0.384 & 0.228 \\
\addlinespace[2pt]
\raisebox{-0.15ex}{
\includegraphics[height=1.5ex]{logos/openai}} & GPT-5.6 Luna & Codex & \texttt{none} & \cellcolor{bandlow}0.099 & \cellcolor{bandlow}0.000 & \cellcolor{bandlow}0.099 & \cellcolor{bandmid}0.114 & \cellcolor{bandlow}0.096 & \cellcolor{bandlow}0.096 & \cellcolor{bandmid}0.137 & \cellcolor{bandlow}0.000 & \cellcolor{bandmid}0.123 & \cellcolor{bandmid}0.149 & 0.091 \\
\raisebox{-0.15ex}{
\includegraphics[height=1.5ex]{logos/openai}} & GPT-5.6 Luna & Codex & \texttt{low} & \cellcolor{bandlow}0.091 & \cellcolor{bandlow}0.091 & \cellcolor{bandlow}0.000 & \cellcolor{bandlow}0.098 & \cellcolor{bandmid}0.175 & \cellcolor{bandlow}0.098 & \cellcolor{bandlow}0.000 & \cellcolor{bandlow}0.000 & \cellcolor{bandmid}0.123 & \cellcolor{bandmid}0.110 & 0.079 \\
\raisebox{-0.15ex}{
\includegraphics[height=1.5ex]{logos/openai}} & GPT-5.6 Luna & Codex & \texttt{medium} & \cellcolor{bandlow}0.098 & \cellcolor{bandlow}0.000 & \cellcolor{bandmid}0.103 & \cellcolor{bandmid}0.132 & \cellcolor{bandmid}0.206 & \cellcolor{bandlow}0.099 & \cellcolor{bandlow}0.084 & \cellcolor{bandlow}0.099 & \cellcolor{bandmid}0.123 & \cellcolor{bandmid}0.297 & 0.124 \\
\raisebox{-0.15ex}{
\includegraphics[height=1.5ex]{logos/openai}} & GPT-5.6 Luna & Codex & \texttt{high} & \cellcolor{bandlow}0.096 & \cellcolor{bandlow}0.073 & \cellcolor{bandlow}0.097 & \cellcolor{bandlow}0.096 & \cellcolor{bandlow}0.095 & \cellcolor{bandlow}0.099 & \cellcolor{bandmid}0.177 & \cellcolor{bandlow}0.098 & \cellcolor{bandmid}0.123 & \cellcolor{bandmid}0.237 & 0.119 \\
\raisebox{-0.15ex}{
\includegraphics[height=1.5ex]{logos/openai}} & GPT-5.6 Luna & Codex & \texttt{xhigh} & \cellcolor{bandlow}0.099 & \cellcolor{bandlow}0.000 & \cellcolor{bandlow}0.098 & \cellcolor{bandmid}0.118 & \cellcolor{bandmid}0.319 & \cellcolor{bandlow}0.100 & \cellcolor{bandmid}0.152 & \cellcolor{bandlow}0.096 & \cellcolor{bandmid}0.117 & \cellcolor{bandmid}0.370 & 0.147 \\
\raisebox{-0.15ex}{
\includegraphics[height=1.5ex]{logos/openai}} & GPT-5.6 Luna & Codex & \texttt{max} & \cellcolor{bandlow}0.091 & \cellcolor{bandlow}0.000 & \cellcolor{bandlow}0.099 & \cellcolor{bandlow}0.098 & \cellcolor{bandmid}0.270 & \cellcolor{bandmid}0.136 & \cellcolor{bandmid}0.143 & \cellcolor{bandmid}0.107 & \cellcolor{bandmid}0.123 & \cellcolor{bandmid}0.326 & 0.139 \\
\bottomrule
\end{tabular}
\end{adjustbox}
\end{table}


\begin{figure}[t]
  \centering
  \includegraphics[width=\linewidth]{fig_system_effort}
  \caption{\textbf{Mean score by system and reasoning effort.} One group per model, one bar per effort
  level, each the mean of that configuration's ten task scores; deeper colour is more effort. The dashed
  line at $0.1$ is the algorithm each repository already ships. Kimi~K3 was run at a single level. Every
  system in the study sits inside the lowest fifth of the scale, and no system rises monotonically with
  effort.}
  \label{fig:bysystem}
\end{figure}

\begin{figure}[t]
  \centering
  \includegraphics[width=\linewidth]{fig_per_task_mapped}
  \caption{\textbf{Mapped score against exploration spend, one panel per task.} Each point is one of
  the 29 configurations; $x$ is the number of output tokens it generated during its four-hour
  exploration of that task (log scale), $y$ is its score after the mapping of
  Table~\ref{tab:mapped}, and the dashed line is $0.1$, the score of the repository's own
  shipped algorithm. We read spend in output tokens rather than total tokens because input counts are
  inflated by each harness's context replay and are not comparable across systems, while output
  tokens are what the reasoning-effort setting actually moves. Each panel keeps its own $y$ range, so
  that the columns sitting on the $0.1$ line stay readable; that line is always within range. Missing
  points are configurations that returned no scorable model.}
  \label{fig:per_task_mapped}
\end{figure}


% Two pieces were parked when this section was cleared, and are held in
% Sections_zh/_parked.tex until it is decided where they belong: the per-task panel
% figure (fig_per_task, still generated by scripts/make_per_task_figure.py) and the
% subsection "What a submission changes", which defines the eight patch families
% behind the introduction's three headline claims.
